\section{The Propers: Interpretive Possibilities}

\begin{studybox}{What this section is}
Everything below is exploratory editorial and AI proposal. None of it is
sourced historical intent, attributed patristic or Doctoral teaching, or
liturgical rubric. Each proposal was checked against the corpus this guide
searched and against a targeted precedent search recorded in
\texttt{research/scope.md}; where no precedent was located, that result is
bounded by the search actually run and is correctable. Each ends with the
strongest objection the author could find against it.
\end{studybox}

\begin{proposal}{1. The congregation signs the first line of its own audit
\normalfont(\emph{Int.} + \emph{Gosp.})}
\textbf{Anchors.} Introit \latin{Suscépimus, Deus, misericórdiam tuam}
(no.\ 1512); Gospel \latin{redde ratiónem villicatiónis tu\ae} (no.\ 1517).

\textbf{Mechanism.} The Introit is in the perfect and the first person plural:
the congregation's opening act is to state, on the record, that it has received
something. Half an hour later the Gospel opens with an employer requiring
exactly that kind of statement from a man who has been handling goods that were
not his. The two texts are not verbally related, but they are the same
\emph{speech act} at two ends of one Mass: an acknowledgement of receipt, and a
demand for one. On this reading the liturgy has arranged for the assembly to
make its declaration \emph{before} the summons is read out, so that when the
summons comes the congregation has already conceded the premise on which the
whole parable turns --- that what is in hand came from somewhere else.

\textbf{Fruit.} The parable stops being a story about a rogue and becomes a
document already initialled. It also explains why the Mass can afford to leave
the parable's scandal unresolved: the hearer's own position has been fixed
liturgically before the difficulty arrives, so he does not need the steward to
be exonerated in order to know what to do.

\textbf{What the element-by-element reading misses.} Read seriatim, the Introit
is a devotional pleasantry about mercy and the Gospel a hard case in moral
theology. The order of speech acts is invisible unless the two are set side by
side.

\textbf{Strongest limit.} There is no verbal dependence whatever: \latin{suscipere}
and \latin{reddere rationem} share no root, and the Introit's \latin{nos} is not
the steward's \latin{ego}. No source consulted connects them. If the Introit was
chosen for its temple imagery --- which its shared use at the Purification makes
likely --- then this correspondence is a coincidence of the calendar and not a
composition.
\end{proposal}

\begin{proposal}{2. The antiphon that lacks a body, and the oration that supplies
one \normalfont(\emph{Int.} + \emph{Postcomm.})}
\textbf{Anchors.} Introit \latin{in médio templi tui} (no.\ 1512), which the
same book prints verbatim at no.\ 2107 for the Purification on 2 February;
Postcommunion \latin{reparátio mentis et córporis} (no.\ 1521).

\textbf{Mechanism.} On 2 February the words ``in the midst of thy temple'' have
a body attached: the Child is carried into the sanctuary, and the antiphon
describes what the arms of Simeon are doing. On the Eighth Sunday the book
reuses the antiphon with no narrative, so the phrase names a place with nothing
in it. The formulary then closes with the only oration of the three that asks
for something bodily --- \latin{reparátio mentis et córporis} --- immediately
after the congregation has received a body. The proposal is that the Sunday's
formulary is completed at its end by the thing its opening antiphon leaves
unspecified: the mercy received in the midst of the temple turns out, at the
Postcommunion, to have been received in bodies.

\textbf{Fruit.} It gives the Purification's antiphon a non-arbitrary reason to
be here, without asserting that the compilers intended the link; and it supplies
a concrete referent for \latin{templi} that does not require allegorising the
church building.

\textbf{What the element-by-element reading misses.} Treated alone, the Introit
is a Zion text and the Postcommunion a brief request for interior effect. The
shared use with 2 February is not visible from this page of the missal at all;
it is visible only from page 467.

\textbf{Strongest limit.} The Pauline identification of the body as a temple
(1 Cor 6:19) is \emph{not} appointed in this Mass, so the bridge depends on a
verse the congregation does not hear. And a simpler explanation of the shared
antiphon is available and probably true: \latin{Suscepimus} is a well-known
chant that fits any Mass with a mercy-and-sanctuary theme, and its reuse says
nothing about either formulary's argument.
\end{proposal}

\begin{proposal}{3. The steward's two refusals are the Offertory's two classes
\normalfont(\emph{Gosp.} + \emph{Off.})}
\textbf{Anchors.} Gospel \latin{fódere non váleo, mendicáre erubésco}
(no.\ 1517); Offertory \latin{Pópulum húmilem salvum fácies \dots\ et óculos
superbórum humiliábis} (no.\ 1518).

\textbf{Mechanism.} The steward names exactly two ways down and rejects both.
Digging is the labour of the poor; begging is the posture of the poor. The next
proper the congregation hears sorts the world into precisely two classes --- the
\latin{pópulus húmilis} who are saved and the \latin{óculi superbórum} who are
brought low --- and the steward has just declined membership in the first on
both available counts. On this reading his real disqualification is not
dishonesty at all. He is commended for foresight and condemned by the following
chant for pride, and the two verdicts are about different things.

\textbf{Fruit.} It offers a way of holding Ambrose and Augustine together
without pretending they agree. Ambrose can be right that the steward is not
censured \emph{for what he foresaw}; Augustine can be right that the act is not
for imitation; and the Offertory can locate the actual fault in a third place
that neither disputes --- the refusal to be lowered. It also gives the offertory
procession an argument rather than a mood.

\textbf{What the element-by-element reading misses.} The Offertory is normally
read as a general thanksgiving verse. Its adjacency to a man who has just said,
in effect, ``I am too good to dig and too proud to beg,'' is lost when the two
are treated as separate items.

\textbf{Strongest limit.} Augustine, expounding the Offertory's own verse, reads
\latin{pópulum húmilem} not as the manual labourer but as those ``who confess
their sins,'' and the proud as those who ``seek to establish their own''
righteousness. That is a moral, not a social, division, and it weakens the neat
mapping of \latin{fódere} onto poverty. The digging half of the proposal is an
extension beyond what any checked witness says.
\end{proposal}

\begin{proposal}{4. The debtors write in their own hands; the assembly says the
Creed in its own voice \normalfont(\emph{Gosp.} + the printed \rubric{Credo})}
\textbf{Anchors.} Gospel \latin{Accipe cautiónem tuam: et sede cito, scribe
quinquagínta \dots\ Accipe lítteras tuas, et scribe octogínta} (no.\ 1517); the
direction \rubric{Credo}, printed in this formulary immediately after
\latin{ætérna tabernácula}.

\textbf{Mechanism.} The parable's one piece of stage business is that the
steward does not alter the bonds himself. He hands each debtor his own document
and has him rewrite it in his own hand --- which is what makes the new figure
binding and the debtor complicit. The missal then prints, as the next thing that
happens in the church, the one text of the Mass that every person present
recites in the first person singular. \latin{Credo} is an instrument executed in
one's own voice, immediately after a story whose hinge is an instrument executed
in one's own hand.

\textbf{Fruit.} It gives the Creed a function in this particular Sunday's
argument instead of being liturgical furniture, and it names something true about
both acts: neither can be performed by proxy, and both create an obligation the
speaker cannot afterwards disown.

\textbf{What the element-by-element reading misses.} A ten-element map of the
propers does not include the Creed at all, because the Creed is not proper.
Anything it contributes is therefore invisible to the standard reading.

\textbf{Strongest limit.} This is the weakest anchor in the section, and
deliberately so marked. \rubric{Credo} is prescribed on Sundays by general
rubric, not chosen for this formulary; the same direction appears on the
Seventh, Ninth, and every other Sunday after Pentecost. Any argument that its
placement here is significant must explain why it is not equally significant
everywhere else, and this proposal cannot.
\end{proposal}

\begin{proposal}{5. A withheld condition and a partial mode of knowing
\normalfont(\emph{Ep.} + \emph{Comm.})}
\textbf{Anchors.} Epistle ending at \latin{coherédes autem Christi} (no.\ 1514),
before Paul's \latin{si tamen compatimur, ut et conglorificemur}; Communion
\latin{Gustáte et vidéte} (no.\ 1520).

\textbf{Mechanism.} The Mass does two withholding things. It hands the
congregation an inheritance and stops one clause short of the terms on which the
inheritance is held; and it then offers a mode of knowing that is by design
incomplete --- taste before sight, the sense that reports presence without
reporting extent. The proposal is that these two are the same gesture at
different points in the rite: the day gives the hearer more than he can yet
account for and less than he will eventually be told, and it does so twice.
Under this reading the truncated verse is not an accident of length but of a
piece with a formulary that ends by asking to \emph{feel} an effect rather than
to see a result.

\textbf{Fruit.} It supplies a reason why this Mass, uniquely among the Sundays
that read Romans 8, can end serenely after a Gospel about a reckoning: the
serenity is not obliviousness but a deliberate partiality of knowledge, which
the Postcommunion then names.

\textbf{What the element-by-element reading misses.} Each element alone is
unremarkable --- pericope boundaries are inherited, and \latin{gustáte et
vidéte} is the commonest of Communion antiphons. Only the conjunction suggests a
pattern, and the conjunction is only visible if one knows where Rom 8:17 breaks.

\textbf{Strongest limit and the disconfirming condition.} The proposal is
vulnerable at its foundation, and the disconfirming evidence sits in the same
book. The 1962 missal appoints Rom 8:18--23, the immediate sequel, on the
\emph{Fourth} Sunday after Pentecost. The condition of co-suffering is therefore
not withheld from the liturgical year at all --- only from this Sunday, and only
because a lection had to end somewhere. If the pericope division is inherited
from an earlier lectionary for reasons of length or of the ancient Roman
station, then nothing is being withheld and this proposal collapses into a
coincidence of scissors. The proposal survives, if at all, as a reading of the
formulary's effect on a hearer, never as a claim about its design.
\end{proposal}

\clearpage
